% from ega4-16-fr.tex
\oldpage[IV-4]{18}
les flèches verticales définissant ici les structures d'Algèbre provenant des
secondes projections; d'ailleurs, si $\sigma$ et $\sigma'$ sont les symétries
canoniques correspondant à $f$ et $f'$ (16.3.4), on a
\[
  \nu_n\circ\rho^*(\sigma)=\sigma'\circ\nu_n
\]
qui fait passer d'un des diagrammes précédents à l'autre. On déduit donc de
(16.4.3.1) un homomorphisme canonique de $\mathcal{O}_{X'}$-Algèbres augmentées
\[
  P^n(u):u^*(\mathscr{P}_{X/S}^n)
  =\mathscr{P}_{X/S}^n\otimes_{\mathcal{O}_X}\mathcal{O}_{X'}
  \longrightarrow\mathscr{P}_{X'/S'}^n
  \tag{16.4.3.3}
  \label{IV.16.4.3.3}
\]
et il résulte de (16.4.3.2) que le diagramme
\[
  \xymatrix{
    \mathcal{O}_{X'} \ar[r]^{\operatorname{id}}
      \ar[d]_{u^*(d_{X/S}^n)} &
      \mathcal{O}_{X'} \ar[d]^{d_{X'/S'}^n} \\
    u^*(\mathscr{P}_{X/S}^n) \ar[r]_{P^n(u)} &
      \mathscr{P}_{X'/S'}^n
  }
  \tag{16.4.3.4}
  \label{IV.16.4.3.4}
\]
est commutatif. On en déduit un homomorphisme de $\mathcal{O}_{X'}$-Algèbres
graduées
\[
  \operatorname{Gr}_\bullet(u):u^*(\mathscr{G}r_\bullet(\mathscr{P}_{X/S}))
  \longrightarrow\mathscr{G}r_\bullet(\mathscr{P}_{X'/S'})
  \tag{16.4.3.5}
  \label{IV.16.4.3.5}
\]
et en particulier un homomorphisme de $\mathcal{O}_{X'}$-Modules
\[
  \operatorname{Gr}_1(u):
  \Omega_{X/S}^1\otimes_{\mathcal{O}_X}\mathcal{O}_{X'}
  \longrightarrow\Omega_{X'/S'}^1
  \tag{16.4.3.6}
  \label{IV.16.4.3.6}
\]
donnant lieu à un diagramme commutatif
\[
  \xymatrix{
    \mathcal{O}_{X'} \ar[r]^{\operatorname{id}}
      \ar[d]_{d_{X/S}\otimes 1} &
      \mathcal{O}_{X'} \ar[d]^{d_{X'/S'}} \\
    \Omega_{X/S}^1\otimes_{\mathcal{O}_X}\mathcal{O}_{X'} \ar[r] &
      \Omega_{X'/S'}^1
  }
  \tag{16.4.3.7}
  \label{IV.16.4.3.7}
\]
\end{env}

\begin{env}[16.4.4]
\label{IV.16.4.4}
Lorsque $S=\operatorname{Spec}(A)$, $S'=\operatorname{Spec}(A')$,
$X=\operatorname{Spec}(B)$, $X'=\operatorname{Spec}(B')$ sont affines, de sorte
que l'on a un diagramme commutatif d'homomorphismes d'anneaux
\[
  \xymatrix{
    B \ar[r] & B' \\
    A \ar[u] \ar[r] & A' \ar[u]
  }
\]
l'image de $\mathfrak{I}_{B/A}$ dans $B'\otimes_{A'}B'$ est contenue dans
$\mathfrak{I}_{B'/A'}$, et l'homomorphisme $\nu_n$ correspond à
l'homomorphisme d'anneaux $P_{B/A}^n\to P_{B'/A'}^n$ déduit de l'homomorphisme
$B\otimes_A B\to B'\otimes_{A'}B'$ par passage aux quotients. L'homomorphisme
(16.4.3.6) correspond à l'homomorphisme défini dans (0, 20.5.4.1), et le
diagramme commutatif (16.4.3.7) au diagramme (0, 20.5.4.2).
\end{env}

